% !TeX root = ../main.tex

\part{Prof. Cincotti}

\section{EEG}

\subsection{Introduction to EEG}

The electroencephalogram, or EEG, is a non-invasive technique used to record electrical potential differences from the scalp. These potentials reflect, in an indirect and spatially filtered way, the activity of large populations of neurons, especially cortical neurons.

Historically, non-invasive human EEG became important in the 1920s and 1930s thanks to Hans Berger, who observed a large and persistent oscillation around 10 Hz, later called the \textit{alpha rhythm}. Later, Adrian and Matthews proposed the idea that some brain rhythms may appear when a cortical area is relatively idle, that is, when it is not strongly engaged in a specific task.

\begin{authornote}
The idea of an ``idling rhythm'' should not be interpreted as the brain being switched off. Rather, it means that an area may show a more synchronized oscillatory pattern when it is not actively processing a specific stimulus or task. When that area becomes functionally engaged, the rhythm often decreases in power.
\end{authornote}

\subsection{Spontaneous Brain Rhythms}

EEG contains spontaneous oscillatory activity. The main frequency bands are usually classified as follows:

\begin{center}
\begin{tabular}{lll}
\toprule
\textbf{Band} & \textbf{Frequency range} & \textbf{General meaning} \\
\midrule
Delta & $< 3.5$ Hz & Slow activity \\
Theta & $4$--$7.5$ Hz & Slow/intermediate activity \\
Alpha & $8$--$13$ Hz & Around 10 Hz rhythm \\
Beta & $14$--$30$ Hz & Faster activity \\
Gamma & $> 30$ Hz & High-frequency activity \\
\bottomrule
\end{tabular}
\end{center}

The alpha rhythm is typically recorded over occipital regions and is strongly related to the visual system. It is often prominent when the eyes are closed and decreases when visual processing increases.

The mu rhythm lies in the same frequency band as alpha, approximately $8$--$13$ Hz, but it has a different functional and spatial meaning. It is recorded mainly over central sensorimotor regions and is related to motor activity, motor preparation, and motor imagery. Its waveform is often more arc-shaped than sinusoidal.

\begin{authornote}
A very important point is that frequency alone is not enough to identify a rhythm. Alpha and mu can both be around 10 Hz. To distinguish them, one must consider location and functional modulation. A 10 Hz rhythm over occipital electrodes that changes with eyes open or closed is likely alpha. A 10 Hz rhythm over central electrodes that changes with movement or motor imagery is likely mu.
\end{authornote}

\subsection{Evoked and Event-Related Potentials}

Besides spontaneous activity, EEG can show responses that are time-locked to external or internal events. A response directly driven by a sensory stimulus is often called an \textit{evoked potential} (EP). The more general term \textit{event-related potential} (ERP) refers to EEG changes related to external stimuli, motor actions, cognitive events, or task-related processes.

ERP analysis became possible with the use of computers, because EEG responses to single events are usually very small and hidden inside the spontaneous EEG. By repeating the same event many times and averaging the EEG epochs aligned to that event, the signal-to-noise ratio improves.

If the EEG response is similar across trials and the background EEG behaves like random noise, averaging $N$ trials improves the signal-to-noise ratio approximately by:

\[
\sqrt{N}.
\]

\begin{authornote}
The key idea is simple: the event-related response is consistent across trials, while the unrelated EEG activity varies more randomly. When we average many trials, the consistent part remains and the random part is reduced.
\end{authornote}

\subsection{EEG Electrodes and Transducers}

EEG electrodes are transducers that interface the biological tissue with the electronic recording system. In the body, electrical currents are mainly carried by ions, whereas in wires and amplifiers they are carried by electrons. The electrode is therefore the interface between ionic conduction and electronic conduction.

Common EEG electrodes are made of silver/silver-chloride, Ag/AgCl. These electrodes are useful because they are relatively non-polarizable and allow stable recordings, including slowly changing potentials.

In EEG terminology, electrode impedance usually refers to the quality of contact between the electrode and the scalp, through conductive gel or another conductive medium.

\begin{authornote}
A useful way to remember this is: the electrode does not ``create'' the EEG signal. It allows the recording system to access potential differences generated by biological activity.
\end{authornote}

\subsection{EEG Amplifiers}

EEG signals are very small, usually in the range of microvolts. Therefore, EEG amplifiers must satisfy strict requirements:

\begin{itemize}[noitemsep]
    \item low noise;
    \item high gain;
    \item high input impedance;
    \item high common-mode rejection ratio (CMRR).
\end{itemize}

EEG amplifiers are typically differential amplifiers. Ideally, they amplify the difference between two input potentials:

\[
V_{\text{out}} \approx A(V_{+} - V_{-}),
\]

where $A$ is the gain of the amplifier.

A ground electrode is also attached to the subject. It provides a common electrical reference for the amplifier system. Its position can be chosen rather freely, provided that it does not introduce a large common-mode signal.

\subsection{Common-Mode Rejection Ratio}

The common-mode rejection ratio, CMRR, measures how well an amplifier rejects signals that are common to both inputs. A good EEG amplifier should amplify the differential signal of interest while suppressing common-mode disturbances such as power-line noise.

The CMRR is defined as the ratio between differential gain and common-mode gain:

\[
\text{CMRR} =
20 \log_{10}
\left(
\frac{\text{differential gain}}
{\text{common-mode gain}}
\right).
\]

High CMRR values, often around 100 dB, indicate better rejection of common-mode noise.

\begin{authornote}
The amplifier does not want to amplify everything that reaches the electrodes. It should amplify what is different between the two inputs and suppress what is common to both. This works well only if the recording system is balanced.
\end{authornote}

\subsection{Input Impedance and Electrode-Scalp Impedance}

A crucial distinction must be made between:

\begin{itemize}[noitemsep]
    \item the input impedance of the amplifier;
    \item the contact impedance between electrode and scalp.
\end{itemize}

The input impedance of the amplifier should be very high, typically in the order of hundreds of megaohms. This prevents the amplifier from loading or attenuating the biological signal.

The electrode-scalp impedance should instead be low and similar across electrodes. Traditional EEG guidelines suggest keeping electrode impedance below approximately $5\,\text{k}\Omega$, although modern high-input-impedance amplifiers allow this requirement to be somewhat relaxed.

If electrode impedances are too high or strongly unbalanced, two problems may occur:

\begin{itemize}[noitemsep]
    \item the biological potential may be attenuated;
    \item the common-mode rejection of the whole recording system may worsen.
\end{itemize}

\begin{noexam}
The green-background slides explain this using the voltage-divider model. The intuition is that if the contact impedance becomes comparable to the amplifier input impedance, part of the biological signal is lost across the contact impedance instead of being measured by the amplifier.

Moreover, if the two electrodes have very different impedances, common-mode noise no longer enters the two amplifier inputs in the same way. This asymmetry reduces the effectiveness of CMRR.
\end{noexam}

\subsection{Standard Electrode Positions: The 10--20 System}

The international 10--20 system defines standard EEG electrode positions using anatomical landmarks:

\begin{itemize}[noitemsep]
    \item nasion;
    \item inion;
    \item left and right preauricular points;
    \item vertex.
\end{itemize}

The name 10--20 comes from the fact that electrode distances are defined as 10\% or 20\% of the distances between reference points on the skull.

Electrode labels indicate scalp regions:

\begin{center}
\begin{tabular}{ll}
\toprule
\textbf{Label} & \textbf{Meaning} \\
\midrule
F & Frontal \\
C & Central \\
P & Parietal \\
O & Occipital \\
T & Temporal \\
z & Midline \\
\bottomrule
\end{tabular}
\end{center}

Odd numbers indicate the left hemisphere, while even numbers indicate the right hemisphere. For example, C3 is a central left electrode, C4 is a central right electrode, and Cz is located on the midline near the vertex.

\begin{authornote}
For motor imagery and mu rhythm studies, C3, Cz, and C4 are especially important because they lie over sensorimotor areas.
\end{authornote}

\subsection{Reference Electrode and Re-Referencing}

Each EEG channel is a potential difference. Therefore, EEG always depends on the chosen reference. Common reference positions include earlobes, mastoids, and the nose. However, these sites are not truly inactive.

The position of the reference electrode can strongly influence the shape and amplitude of EEG waveforms. For this reason, reference choice is not a minor technical detail.

One common approach is the common average reference (CAR), where the instantaneous average of all channels is subtracted from each channel. In ideal theoretical conditions, this would approximate a reference at infinity. In practice, these ideal conditions are never perfectly satisfied.

\begin{authornote}
A useful sentence to remember is: the EEG reference changes what we see. The same neural activity may appear with different amplitude or waveform depending on the reference used.
\end{authornote}

\subsection{Monopolar and Bipolar Montages}

In a monopolar montage, all channels are measured with respect to a common reference. This is useful for studying spatial potential distributions over the scalp.

In a bipolar montage, each channel is the difference between two nearby electrodes. Bipolar montages are widely used in clinical EEG because they help highlight local changes and phase reversals.

\begin{authornote}
The monopolar montage answers the question: ``What is the potential at each electrode relative to a reference?'' The bipolar montage answers: ``How does the potential change between nearby electrodes?''
\end{authornote}

\subsection{Good Practices in EEG Data Collection}

Good EEG data require careful experimental preparation. Recordings should be performed with alert and cooperative subjects in conditions as free as possible from artifacts and noise.

The recording session should be divided into smaller runs or blocks, separated by breaks. Subjects should be reminded to:

\begin{itemize}[noitemsep]
    \item relax their muscles;
    \item minimize blinking;
    \item avoid head and body movements;
    \item avoid jaw clenching, swallowing, and facial tension.
\end{itemize}

The experimenter should continuously monitor the EEG online to detect artifacts or technical problems. Accurate logging is also essential: each run should be documented, and relevant metadata or time markers should be saved with the data.

\begin{authornote}
EEG cannot be fully fixed after acquisition. Signal processing can reduce artifacts, but bad acquisition, missing markers, or unstable electrodes can seriously compromise the analysis.
\end{authornote}

\subsection{EEG Artifacts}

An EEG artifact is any potential difference due to an extracerebral source. Artifacts may be biological or technical.

Biological artifacts include:

\begin{itemize}[noitemsep]
    \item eye movements and blinks;
    \item muscle activity;
    \item cardiac activity;
    \item sweating;
    \item body and head movements.
\end{itemize}

Technical artifacts include:

\begin{itemize}[noitemsep]
    \item power-line noise;
    \item bad electrode contact;
    \item electrode detachment;
    \item electrical interference from external devices;
    \item ADC saturation.
\end{itemize}

Artifact prevention is preferable to post-hoc artifact removal.

\begin{authornote}
Many artifacts have amplitudes much larger than the actual EEG signal. Once they contaminate the recording, they can sometimes be attenuated, but this may also remove useful neural information.
\end{authornote}

\subsubsection{Eye-Related Artifacts}

The eye behaves like an electrical dipole: the cornea is more positive than the retina. When the eye moves, this dipole moves inside the volume conductor and produces large potentials detectable by EEG electrodes.

Eye movement artifacts can be around $0.5\,\text{mV}$, much larger than ordinary EEG rhythms. Eye blinks produce large, slow, often monophasic deflections lasting around 200--400 ms. These artifacts are especially visible over frontal electrodes.

\begin{authornote}
If very large slow waves appear mainly over frontal channels, the first suspect should be blinks or eye movements, not necessarily brain activity.
\end{authornote}

\subsubsection{Muscle Artifacts}

Muscle activity, EMG, can contaminate EEG recordings. Muscle artifacts may range from hundreds of microvolts to millivolts and have a broad frequency spectrum, from tens of Hz to several kHz.

This is especially problematic because EMG overlaps with beta and gamma EEG bands. Common sources include jaw clenching, swallowing, tongue tension, forehead contraction, and neck muscle activity.

\begin{authornote}
If the EEG trace looks very fast, irregular, and high-frequency, especially over temporal or frontal areas, EMG contamination should be suspected.
\end{authornote}

\subsubsection{Cardiac Artifacts}

The heart generates electrical potentials recorded as ECG or EKG. These can contaminate EEG, especially when the reference electrode is not placed on the head.

Another cardiac-related artifact is the ballistocardiogram, caused by small electrode movements due to nearby blood vessels pulsating. This artifact can sometimes be reduced by slightly moving the electrode.

\subsubsection{Sweating and Movement Artifacts}

Sweating produces electrodermal responses, which are slow and high-amplitude potentials, typically below $0.5$ Hz.

Head and body movements may also produce artifacts. These may come from muscle activity or from mechanical displacement of the charged double layer at the electrode-skin interface.

\begin{authornote}
Slow artifacts are not harmless. They can strongly affect baseline correction and ERP analysis, especially when studying slow components.
\end{authornote}

\subsection{Power-Line Noise}

Power-line noise arises from capacitive coupling with alternating current from electrical power supplies. In Europe, it is typically centered around 50 Hz; in other countries, it may be 60 Hz. Harmonics may also appear.

A notch filter can selectively attenuate 50 Hz noise. However, it should be used carefully if the frequency band of interest is close to the rejected frequency.

\begin{authornote}
For mu and beta analysis, typically 8--30 Hz, a 50 Hz notch filter is usually not too problematic. For gamma-band analysis, however, it may remove or distort useful information.
\end{authornote}

\subsection{Offline EEG Inspection}

Before any advanced processing, raw EEG data should be visually inspected. Bad channels or bad epochs must be identified and removed when necessary.

The principle is:

\[
\text{garbage in, garbage out}.
\]

However, one should avoid rejecting too much data. The decision depends on the specific analysis and on which artifacts are relevant for the research question.

\begin{authornote}
There is no universal cleaning pipeline. If the goal is to study frontal ERPs, blinks are a major problem. If the goal is to study occipital alpha, isolated frontal blinks may be less damaging, but baseline shifts and movement artifacts may still be critical.
\end{authornote}

\subsection{Segmentation and Averaging}

For ERP or EP analysis, the continuous EEG recording is segmented into epochs. Each epoch usually includes:

\begin{itemize}[noitemsep]
    \item a pre-stimulus or pre-event baseline period;
    \item a post-stimulus or post-event period.
\end{itemize}

The epochs are aligned to the event of interest and averaged sample by sample. This synchronized averaging preserves the event-related response and reduces unrelated spontaneous EEG activity.

ERP peaks are usually characterized by:

\begin{itemize}[noitemsep]
    \item amplitude, measured relative to the pre-stimulus baseline;
    \item latency, measured relative to the stimulus or event onset.
\end{itemize}

ERP components are often named using polarity and nominal latency. For example:

\begin{itemize}[noitemsep]
    \item P100: positive peak around 100 ms;
    \item N170: negative peak around 170 ms;
    \item P300: positive peak around 300 ms.
\end{itemize}

\begin{authornote}
The number in the component name is nominal, not exact. A negative peak at 108 ms may still be called N100 if it corresponds to the physiological N100 component.
\end{authornote}

\subsection{Stimulus Timing: SOA, ISI, and ITI}

In EEG experiments, timing is crucial. The main timing terms are:

\begin{center}
\begin{tabular}{ll}
\toprule
\textbf{Term} & \textbf{Meaning} \\
\midrule
SOA & Time between the onset of one stimulus and the onset of the next \\
ISI & Time between the end of one stimulus and the onset of the next \\
ITI & Time between the beginning of one trial and the beginning of the next \\
\bottomrule
\end{tabular}
\end{center}

If the stimulus duration is long, SOA and ISI can be very different. ERP amplitude and shape may depend strongly on stimulus timing. Long-latency components are especially sensitive to fast stimulus repetition.

\begin{authornote}
There is a trade-off: more trials improve the signal-to-noise ratio, but if stimuli are too close in time, responses may overlap or decrease in amplitude.
\end{authornote}

\subsection{Evoked and Induced Activity}

Evoked activity is both time-locked and phase-locked to the stimulus. Therefore, it can be revealed by conventional time-locked averaging.

Induced activity is related to the event but not phase-locked. It may occur at a similar frequency after each stimulus, but with variable phase or latency. Because of this jitter, it may disappear in the average waveform.

To study induced activity, one can analyze amplitude or power in a specific frequency band, for example by filtering the signal and then rectifying or squaring it before averaging.

\begin{authornote}
ERP analysis is good for activity that is aligned in time and phase. If the brain responds with a change in oscillatory power but with variable phase, time-frequency analysis or ERD/ERS analysis is more appropriate.
\end{authornote}

\subsection{Event-Related Desynchronization and Synchronization}

Event-related desynchronization and synchronization, ERD/ERS, quantify changes in EEG power relative to a baseline period.

The typical procedure is:

\begin{enumerate}[noitemsep]
    \item band-pass filter the EEG in the frequency band of interest;
    \item square the signal to estimate power;
    \item average across trials;
    \item smooth the result using short time windows;
    \item compute the percentage change relative to baseline.
\end{enumerate}

The general formula is:

\[
\text{ERD/ERS}(\%) =
\frac{P_{\text{task}} - P_{\text{baseline}}}
{P_{\text{baseline}}}
\cdot 100.
\]

If the value is negative, there is ERD. If the value is positive, there is ERS.

\begin{itemize}[noitemsep]
    \item \textbf{ERD}: decrease in power relative to baseline.
    \item \textbf{ERS}: increase in power relative to baseline.
\end{itemize}

In motor imagery, ERD in the mu band over C3 during right-hand imagery suggests involvement of the left sensorimotor cortex.

\begin{authornote}
ERD does not simply mean ``less brain activity''. In sensorimotor EEG, ERD often means that the synchronized rhythm is disrupted because the cortical area is functionally engaged.
\end{authornote}

\subsection{Time-Frequency Analysis}

Time-frequency analysis is used to observe how the spectral content of EEG changes over time. It is especially useful for studying induced activity and ERD/ERS.

There is always a trade-off between time resolution and frequency resolution. Short time windows provide good temporal resolution but poor frequency resolution. Long time windows provide good frequency resolution but poor temporal resolution.

\begin{authornote}
This trade-off is one of the main practical limitations of time-frequency EEG analysis. You cannot arbitrarily obtain perfect timing and perfect frequency precision at the same time.
\end{authornote}

\subsection{Live EEG Demonstration}

The live EEG demonstration used a wireless 32-channel amplifier and active Ag/AgCl electrodes. These electrodes include on-board impedance converters, producing a low-impedance signal through the wire and therefore reducing susceptibility to interference.

The setup requires symmetric cap placement based on anatomical landmarks such as nasion, inion, preauricular points, and vertex. Impedances must be monitored. A scientific target is usually below $5\,\text{k}\Omega$, while a more relaxed demonstration target may be below $25\,\text{k}\Omega$.

The demonstration also illustrates the practical effect of filters:

\begin{itemize}[noitemsep]
    \item high-pass filters reduce slow drifts;
    \item low-pass filters reduce high-frequency components;
    \item notch filters attenuate power-line noise at 50 Hz.
\end{itemize}

\begin{authornote}
The demonstration connects the theoretical concepts to practical observation. Removing the notch filter may reveal 50 Hz noise; asking the subject to clench the jaw may reveal EMG; changing the high-pass filter may strongly affect slow baseline fluctuations.
\end{authornote}

\subsection{Core Exam Takeaways}

The most important points to remember are:

\begin{enumerate}
    \item EEG measures potential differences, not absolute brain activity.
    \item Alpha and mu rhythms may have similar frequencies but different locations and meanings.
    \item ERP analysis relies on time-locked averaging to improve signal-to-noise ratio.
    \item Evoked activity is phase-locked; induced activity is not necessarily phase-locked.
    \item Amplifiers need high input impedance, while electrode-scalp impedance should be low and balanced.
    \item High CMRR is essential to reject common-mode disturbances such as power-line noise.
    \item The reference electrode is not neutral and can influence EEG waveform shape and amplitude.
    \item Artifacts can be biological or technical and are often larger than the brain signal.
    \item ERD/ERS measure percentage changes in band power relative to baseline.
    \item Good data acquisition is more important than aggressive post-processing.
\end{enumerate}

\subsection{True Statements}

This section collects the true statements explicitly provided in Prof. Cincotti's EEG lectures. They are useful as compact exam-oriented facts.

\subsubsection{Spontaneous EEG Rhythms and Event-Related Responses}

\begin{enumerate}[label=\arabic*.]
    \item The alpha rhythm is an oscillatory component of the spontaneous EEG.
    
    \item The frequency of oscillation of the alpha rhythm is around 10 Hz.
    
    \item The amplitude of the EEG in the alpha band decreases when the generating region of the cerebral cortex becomes engaged in a functional task, such as a visual or motor task.
    
    \item The proper visual alpha rhythm is generated in the occipital lobe of the cerebral cortex.
    
    \item The mu rhythm is an oscillatory component of the spontaneous EEG whose frequency lies in the alpha band.
    
    \item The mu rhythm is generated in the central regions of the cerebral cortex.
    
    \item The oscillations of the mu rhythm are more arc-shaped rather than resembling a regular sine wave.
    
    \item The amplitude of the alpha rhythm is not constant, but rather waxes and wanes with irregular periods, with changes often occurring after intervals in the order of one second.
    
    \item The alpha band of the EEG is conventionally limited between 8 and 13 Hz.
    
    \item The delta and theta frequency bands identify frequencies lower than those in the alpha band.
    
    \item The beta and gamma frequency bands identify frequencies higher than those in the alpha band.
    
    \item The alpha rhythm can be observed by filtering the spontaneous EEG signal using a narrow-band filter, with cutoff frequencies approximately at 8 and 13 Hz.
    
    \item Evoked potentials are deflections of the EEG signal following the presentation of a sensory input.
    
    \item Event-related potentials include evoked potentials, as well as EEG responses to motor or cognitive events.
\end{enumerate}

\subsubsection{Electrodes, Impedance, Amplifiers, and Montages}

\begin{enumerate}[label=\arabic*.]
    \item EEG electrodes made of silver with a layer of silver chloride, Ag/AgCl, are non-polarizable and therefore allow the recording of extremely slow-changing potentials.
    
    \item In EEG terminology, impedance is a measure of the quality of the contact between electrode and scalp, through the conductive gel.
    
    \item Contact impedance of the electrodes is measured in kilo-ohms, $\text{k}\Omega$, and must be measured using an alternating current.
    
    \item Gold electrodes are polarizable; therefore, they should not be used to measure slowly changing potentials.
    
    \item The Common-Mode Rejection Ratio, CMRR, of a differential amplifier characterizes the ratio between the gain of the difference of potential between input electrodes and the gain of the common potential with respect to ground.
    
    \item The CMRR is usually expressed in decibels, dB, and high values characterize better amplifiers.
    
    \item The advantage of a high-CMRR amplifier is that it suppresses common-mode disturbances, such as power-line noise at 50 Hz.
    
    \item The measurement of a single EEG signal requires three electrodes: two electrodes as inputs to the differential amplifier and one electrode to provide the ground potential.
    
    \item The input impedance of a biosignal amplifier must be many orders of magnitude higher than the contact impedance of the electrodes. It is usual to have input impedances in the order of $10^8\,\Omega$.
    
    \item EEG guidelines suggest keeping the electrodes' contact impedance below $5\,\text{k}\Omega$. The use of modern amplifiers with high input impedance allows this requirement to be relaxed.
    
    \item One reason for keeping the electrodes' contact impedance much lower than the amplifier's input impedance is that the resulting voltage divider would otherwise reduce the measured potential.
    
    \item The difference between contact impedances of electrodes should be small compared to the input difference of the differential amplifier; otherwise, the resulting unbalance compromises its common-mode rejection capability.
    
    \item The electrode labels of the International 10--20 System use the first letter of the four lobes of the cerebral cortex, Frontal, Parietal, Occipital, Temporal, plus ``C'' for central electrodes over the central sulcus.
    
    \item In the electrode labels of the International 10--20 System, odd numbers designate electrodes on the left side and even numbers designate electrodes on the right side.
    
    \item The International 10--20 System for EEG electrode placement takes its name from the fact that the distance between adjacent electrodes is 10\% or 20\% of the distance between pairs of reference points on the skull, such as nasion and inion or the preauricular points.
    
    \item In monopolar EEG recordings, a single reference potential is used for all recorded channels. The reference electrode is placed on scalp positions assumed to be far from the electrical sources of interest, such as the earlobes.
    
    \item The position of the reference electrode can strongly influence the shape and amplitude of EEG potentials. The profile of scalp topographies, disregarding the actual potential value, is not influenced.
    
    \item In bipolar EEG recordings, each channel is the difference of potential between two adjacent electrodes. This montage is mostly used in clinical settings to highlight features of interest during visual inspection of the waveforms.
    
    \item EEG signals recorded in monopolar configuration can be re-referenced to the Common Average Reference, CAR, by subtracting from each channel the instantaneous average of all channels. In ideal conditions, this would approximate taking the reference potential at infinity.
\end{enumerate}

\subsubsection{EEG Artifacts}

\begin{enumerate}[label=\arabic*.]
    \item An EEG artifact is a potential difference due to sources outside the brain.
    
    \item Artifacts can have biological origin, such as eyes, muscles, and heart; they can be due to external electromagnetic generators, such as power supply or engines; or they can be due to events affecting the recording setup, such as electrode movements, loss of contact, or saturation of the ADC.
    
    \item Artifacts can be partially attenuated or removed through signal processing during data analysis, but it is always preferable to make all efforts to prevent them during signal acquisition.
    
    \item The eye is more positive in its frontal part, the cornea, than in its posterior part, the retina. Therefore, eye movements can generate large EOG artifacts on the EEG, of the order of $500\,\mu\text{V}$.
    
    \item During an eyeblink, the eyelid shortcuts the positive potential of the external surface of the eye, causing positive EEG deflections with amplitudes of several hundreds of microvolts and durations of a few hundreds of milliseconds.
    
    \item A sudden upward or downward movement of the eyes generates a positive or negative deflection of EEG potentials, respectively, on frontal EEG channels.
    
    \item A sudden movement of the eyes to the right or left generates a positive or negative deflection of EEG potentials, respectively, on the EEG channel F8.
    
    \item The electrical activity of muscles, EMG, has spectral content starting at frequencies of about 20 Hz and above, thus affecting the beta and gamma bands of the EEG signal.
    
    \item The amplitude of the EMG generated by muscles in the head can be ten times higher than the spontaneous EEG signal, thus in the order of $1\,\text{mV}$.
    
    \item EMG artifacts can easily appear in EEG recordings unless subjects are specifically instructed by the experimenter on how to relax their face, tongue, and throat muscles.
    
    \item Heart activity can contaminate EEG recordings because an electrocardiographic artifact, ECG or EKG, can directly affect the potentials, especially if the reference electrode is not placed on the head.
    
    \item Heart activity can also contaminate EEG recordings because a ballistocardiographic artifact is indirectly generated by the pulse of a blood vessel causing movements of a nearby electrode.
    
    \item Sweating can affect EEG, causing slow-changing and high-amplitude artifacts below 0.5 Hz, up to a few millivolts.
    
    \item Power-line noise is an artifact caused by capacitive coupling between conductors carrying alternating current, typically at 50 Hz, and the recording setup, including the subject.
    
    \item Power-line noise affects a very narrow frequency band of the recorded signal around 50 Hz or 60 Hz, depending on the power-line frequency. Other frequency bands can be affected at multiples, typically odd multiples, of 50 or 60 Hz.
    
    \item Power-line noise is accentuated by asymmetries in the recording electrode pairs, such as impedance differences and cable path differences, because asymmetries prevent the noise from being rejected by the amplifier's common-mode rejection capabilities.
    
    \item Notch filters effectively remove power-line noise because they selectively reject the narrow band affected by the artifact, preserving almost entirely the useful signal.
    
    \item Movement of the subject's head may produce slow artifacts on the EEG recording, whose waveform is closely related to the time course of the movement.
    
    \item Since movement-related potentials can originate from the mechanical displacement of the charged double layer at the electrode interface, these artifacts are less pronounced when non-polarizable electrodes are used.
\end{enumerate}

\subsubsection{ERP, EP, ERD/ERS, and Time-Frequency Analysis}

\begin{enumerate}[label=\arabic*.]
    \item As a preliminary step to EEG data analysis, the following data can be discarded: one or more channels, if they are extensively contaminated by artifacts; and all time intervals, or epochs, in which artifacts appear. Both strategies can be applied to the same dataset.
    
    \item The estimation of event-related or evoked potentials, ERPs and EPs, requires the acquisition of numerous repetitions, typically tens or hundreds, of the stimulus or event that evoked or induced the potential.
    
    \item When recording ERPs or EPs, whose peak amplitudes range from a few microvolts down to a fraction of a microvolt, the spontaneous EEG, whose amplitude is of tens of microvolts, is to be considered noise that completely masks the EPs or ERPs on the recorded waveform.
    
    \item In ERP analysis, the averaging procedure consists of segmenting epochs, or trials, with fixed duration from the raw recording, each aligned to a repetition of the event, and then performing a synchronized average, that is, averaging all corresponding samples sharing the same latency across the set of trials.
    
    \item Synchronized averaging of $N$ trials preserves the amplitude of the ERP and reduces the amplitude of the background spontaneous EEG activity by a factor $\sqrt{N}$, under commonly verified hypotheses.
    
    \item The amplitude of ERPs is measured with respect to a baseline epoch, usually preceding the stimulus, in which the amplitude is assumed to be zero.
    
    \item The latency of an ERP peak is measured with respect to the relevant event, usually the presentation of a stimulus, rather than with respect to the beginning of the waveform.
    
    \item In an ERP, peaks are named with a leading P if the peak has positive polarity and with a leading N if the peak has negative polarity.
    
    \item In an ERP, peaks are often named with a trailing number indicating the nominal latency in milliseconds. In older conventions, the trailing number represents the order of the peak within the ERP.
    
    \item A negative peak in an ERP recorded on a specific subject with a latency of 108 ms may still be named N100 if it matches the physiological phenomenon of the nominal N100 component.
    
    \item The Stimulus Onset Asynchrony, SOA, measures the time interval between the onsets of two successive stimuli in a train. If each trial contains only one stimulus, it is equivalent to the Inter-Trial Interval, ITI.
    
    \item The Inter-Stimulus Interval, ISI, measures the time interval between the end of one stimulus and the beginning of the following one. It equals the SOA minus the stimulus duration.
    
    \item In an ERP, the response to a stimulus has reduced amplitude when the SOA is too short. Long-latency components of the ERP are especially sensitive to this decrease.
    
    \item Brain activity in response to a stimulus can be phase-locked to the event, meaning that the whole time course, including positive and negative peaks, has the same latency in every repetition. This activity is called evoked.
    
    \item Brain activity in response to a stimulus can be non-phase-locked, meaning that it shows variable latency, or jitter, at each repetition. This activity is called induced.
    
    \item The averaging procedure can reliably uncover components of an ERP corresponding to evoked brain activity.
    
    \item Induced activity is often examined by analyzing the envelope of the EEG in a relevant frequency band, that is, by rectifying or squaring the pass-band filtered trials before averaging them.
    
    \item In EEG terminology, synchronization and desynchronization are synonymous with an increase and decrease, respectively, of waveform amplitude and thus power.
    
    \item Event-Related Desynchronization/Synchronization, ERD/S, quantifies changes of EEG power relative to a baseline period, expressed as percentage change. ERD/S is usually evaluated over a whole range of relevant latencies with respect to the event.
    
    \item Computation of ERD/S from a set of $N$ EEG trials requires the following steps: band-pass filtering in the relevant frequency band; squaring each sample; synchronized averaging across trials; smoothing by averaging samples within short adjacent time windows covering the whole time course; and computing the relative percentage change with respect to the average value in a baseline period.
    
    \item An alternative algorithm to evaluate ERD/S uses rectification instead of squaring, thus estimating the amplitude of the signal rather than its power.
    
    \item Time-frequency analysis allows the analysis and visualization of changes in the spectral components of a signal over time.
    
    \item Time resolution and frequency resolution cannot be freely chosen: the higher one is, the lower the other becomes.
\end{enumerate}

\section{Signal Processing}

\subsection{Role of Signal Processing in Biomedical Measurements}

A typical biomedical measurement system starts from physiological or external energy and converts it into a signal that can be processed by a computer. The general pipeline is:

\[
\text{physiological phenomenon}
\rightarrow
\text{transducer}
\rightarrow
\text{amplifier}
\rightarrow
\text{analog processing}
\rightarrow
\text{ADC}
\rightarrow
\text{digital signal processing}.
\]

In neuroengineering, signal processing has two main roles:

\begin{itemize}[noitemsep]
    \item preprocessing or conditioning, that is, improving the quality of the signal before further analysis;
    \item feature extraction, that is, transforming raw signals into meaningful numerical descriptors.
\end{itemize}

Examples of preprocessing include filtering, artifact attenuation, and signal segmentation. Examples of feature extraction include estimating signal power in a frequency band, computing the PSD, extracting ERP amplitudes, or measuring EMG amplitude features.

\begin{authornote}
The key idea is that biomedical signals are rarely used directly in their raw form. We usually process them to remove irrelevant components and extract features that are more directly related to the physiological question.
\end{authornote}

\subsection{Analog-to-Digital Conversion}

Analog-to-Digital Conversion, ADC, transforms a continuous analog signal into a digital sequence of numbers. It consists of two main operations:

\[
\text{ADC} =
\text{sampling}
+
\text{quantization}.
\]

Sampling discretizes time. Quantization discretizes amplitude.

\begin{itemize}[noitemsep]
    \item \textbf{Sampling}: the signal is measured at discrete time instants.
    \item \textbf{Quantization}: each measured amplitude is approximated to the nearest allowed digital level.
\end{itemize}

\begin{authornote}
A useful way to remember the difference is: sampling decides \emph{when} we measure; quantization decides \emph{how precisely} we represent the amplitude.
\end{authornote}

\subsection{Quantization and Quantization Noise}

Quantization approximates each analog value to one of a finite number of digital levels. If the ADC uses $N_{\text{bits}}$ bits, the number of available quantization levels is approximately:

\[
2^{N_{\text{bits}}}.
\]

The width of a quantization interval is called LSB, Least Significant Bit. If the input range of the ADC is $R$, then approximately:

\[
LSB =
\frac{R}{2^{N_{\text{bits}}}}.
\]

Quantization introduces an error, called quantization noise. Under standard assumptions, this noise can be modeled as additive, uniformly distributed, zero-mean noise. Its standard deviation is:

\[
\sigma_q =
\frac{LSB}{\sqrt{12}}.
\]

Increasing the number of bits reduces the quantization interval and therefore reduces the quantization noise.

\begin{authornote}
Quantization noise is not the same as physiological noise. It is an artifact introduced by the digital representation itself. A good ADC should make quantization noise smaller than the intrinsic noise of the input signal.
\end{authornote}

\subsection{Input Range, Precision, and Clipping}

The ADC input range must be chosen carefully.

If the range is too large, the quantization interval becomes larger, increasing quantization noise. If the range is too small, large-amplitude components may exceed the ADC range and produce clipping or saturation.

\[
\text{large range}
\rightarrow
\text{larger quantization error},
\]

\[
\text{small range}
\rightarrow
\text{higher risk of clipping}.
\]

\begin{authornote}
This is a practical trade-off. You want the range to be large enough to avoid saturation, but not unnecessarily large, otherwise you waste amplitude resolution.
\end{authornote}

\subsection{Sampling and the Nyquist-Shannon Theorem}

The Nyquist-Shannon sampling theorem states that a continuous signal can be correctly sampled only if it does not contain frequency components above one half of the sampling frequency.

If the sampling frequency is $f_s$, the Nyquist frequency is:

\[
f_N =
\frac{f_s}{2}.
\]

For proper sampling, the maximum frequency component of the signal must satisfy:

\[
f_{\max} < f_N.
\]

Equivalently:

\[
f_s > 2 f_{\max}.
\]

\begin{authornote}
The theorem does not say that every signal sampled above twice the frequency of interest is automatically usable. It assumes that frequencies above Nyquist have been removed before sampling.
\end{authornote}

\subsection{Aliasing}

Aliasing occurs when a signal contains frequency components above the Nyquist frequency before sampling. These high-frequency components are incorrectly represented as lower-frequency components in the sampled signal.

For a sinusoidal component with frequency:

\[
f_0 \in (f_N, f_s),
\]

the aliased frequency is:

\[
f_{\text{alias}}
=
f_s - f_0.
\]

Aliasing is dangerous because, after sampling, the false low-frequency component cannot be distinguished from a real low-frequency component.

\begin{authornote}
Aliasing is not just a loss of information. It is worse: high-frequency content is actively misrepresented as low-frequency content.
\end{authornote}

\subsection{Anti-Aliasing Filtering}

Aliasing can be prevented by applying an analog low-pass filter before the ADC. This filter is called an anti-aliasing filter.

Its role is to remove frequency components above the Nyquist frequency before sampling:

\[
f_{\text{cutoff}} < f_N.
\]

This must be done in the analog domain, before digitization.

\begin{authornote}
A digital filter applied after sampling cannot undo aliasing, because once the high-frequency component has been folded into the low-frequency range, it is already indistinguishable from genuine low-frequency content.
\end{authornote}

\subsection{Analog Signal Conditioning}

Real ADCs are limited by input range and maximum sampling frequency. Therefore, the analog signal should be conditioned before conversion so that:

\begin{itemize}[noitemsep]
    \item the signal range does not exceed the ADC input range;
    \item the signal bandwidth does not exceed what is needed for faithful representation;
    \item high-amplitude artifacts are attenuated when possible;
    \item frequencies above Nyquist are removed.
\end{itemize}

Analog filtering is therefore a key step before ADC.

\subsection{Statistics, Probability, and Noise}

Signals can be deterministic or stochastic.

A deterministic signal is fully described by an explicit rule or waveform. Examples include sine waves, square waves, and triangular waves.

A stochastic signal is described statistically. Its exact sample values cannot be predicted deterministically, but its statistical properties can be described.

Probability theory models random variables and random processes. Statistics describes empirical observations and estimates the parameters of probabilistic models.

\begin{authornote}
In signal processing, the word ``sample'' can be ambiguous. In statistics, a sample is a set of observations. In signal processing, a sample is one measurement of a signal at a specific time instant.
\end{authornote}

\subsection{Basic Signal Measures}

Some common measures for signal characterization are:

\begin{itemize}[noitemsep]
    \item mean;
    \item variance;
    \item standard deviation;
    \item root mean square, RMS;
    \item average rectified value, ARV.
\end{itemize}

The mean of a signal with samples $x_i$ is:

\[
\bar{x}
=
\frac{1}{N}
\sum_{i=1}^{N}
x_i.
\]

The variance is estimated as:

\[
s^2
=
\frac{1}{N-1}
\sum_{i=1}^{N}
(x_i - \bar{x})^2.
\]

The standard deviation is:

\[
s =
\sqrt{s^2}.
\]

The RMS is:

\[
RMS
=
\sqrt{
\frac{1}{N}
\sum_{i=1}^{N}
x_i^2
}.
\]

The ARV is:

\[
ARV
=
\frac{1}{N}
\sum_{i=1}^{N}
|x_i|.
\]

For zero-mean signals, RMS is closely related to standard deviation.

\begin{authornote}
RMS and ARV are often used when we care about signal amplitude but the signal oscillates around zero. This is particularly important for EMG, where raw positive and negative phases would cancel out if we used the mean.
\end{authornote}

\subsection{White Noise and Gaussian Noise}

White noise is a stochastic signal whose samples are uncorrelated. Knowing one sample does not help predict another sample.

In the frequency domain, ideal white noise has a flat spectrum, meaning it has the same power at all frequencies.

Gaussian noise is noise whose sample amplitude follows a normal distribution, often with zero mean.

\[
x \sim \mathcal{N}(0,\sigma^2).
\]

\begin{authornote}
Whiteness describes correlation across time. Gaussianity describes the distribution of amplitudes. A signal can be white without being Gaussian, and Gaussian without being perfectly white.
\end{authornote}

\subsection{Central Limit Theorem and Averaging}

The Central Limit Theorem states that the average of many independent and identically distributed random variables tends to a Gaussian distribution as the number of variables increases.

If $N$ independent random variables have standard deviation $\sigma$, the standard deviation of their average is:

\[
\sigma_{\text{avg}}
=
\frac{\sigma}{\sqrt{N}}.
\]

This is why averaging reduces random noise.

For example, if $N$ EEG trials contain only spontaneous EEG noise with RMS equal to $\sigma$, their synchronized average has RMS approximately:

\[
RMS_{\text{avg}}
=
\frac{\sigma}{\sqrt{N}}.
\]

\begin{authornote}
This is the mathematical reason why ERP averaging works: the event-related component is time-locked and remains, while unrelated random activity is reduced by averaging.
\end{authornote}

\subsection{Fourier Analysis}

Fourier analysis decomposes a signal into a sum of sine and cosine components. Each sinusoidal component carries power at a specific frequency.

This allows us to move from the time domain:

\[
x(t)
\]

to the frequency domain:

\[
X(f).
\]

The frequency domain representation tells us which frequencies are present in the signal and how much amplitude or power each frequency carries.

\subsection{Discrete Fourier Transform}

The Discrete Fourier Transform, DFT, transforms a time-limited and sampled signal into a discrete frequency-domain representation.

If a signal has $N$ samples, its DFT contains $N$ complex values. For a real-valued time-domain signal, the frequency spectrum is symmetric:

\[
X(-f)
=
X^*(f).
\]

For this reason, only the first half of the spectrum, up to the Nyquist frequency, is needed to characterize the signal.

The DFT represents each frequency component through:

\begin{itemize}[noitemsep]
    \item magnitude, or absolute value;
    \item phase, or angle.
\end{itemize}

\begin{authornote}
Magnitude tells us how much of a frequency is present. Phase tells us where that sinusoidal component is shifted in time.
\end{authornote}

\subsection{Fast Fourier Transform}

The Fast Fourier Transform, FFT, is an efficient implementation of the DFT. It is especially efficient when the number of samples is a power of two.

\[
FFT = \text{efficient algorithm for computing the DFT}.
\]

\subsection{Frequency Resolution}

If a signal has total duration $T_{\max}$, the distance between adjacent DFT frequency bins is:

\[
\Delta f =
\frac{1}{T_{\max}}.
\]

Thus, longer signals provide better frequency resolution.

\[
T_{\max} \uparrow
\quad \Rightarrow \quad
\Delta f \downarrow.
\]

\begin{authornote}
Frequency resolution is determined by observation time, not directly by sampling frequency. Sampling frequency determines the maximum representable frequency; signal duration determines the spacing between frequency bins.
\end{authornote}

\subsection{Zero-Padding}

Zero-padding means adding samples with zero amplitude at the end of a signal before computing the DFT.

Zero-padding reduces the distance between adjacent displayed frequency samples, making the spectrum look smoother.

However, it does not add new information and does not truly improve the physical frequency resolution determined by the original signal duration.

\begin{authornote}
Zero-padding is interpolation in the frequency domain. It helps visualization and peak localization, but it does not create new data.
\end{authornote}

\subsection{Windowing and Spectral Leakage}

When we analyze a finite section of a longer signal, we are effectively multiplying the signal by a windowing function.

Using a rectangular window may introduce discontinuities at the edges of the selected segment. These discontinuities cause spectral leakage: energy from one frequency spreads into nearby frequencies.

Spectral leakage is characterized by:

\begin{itemize}[noitemsep]
    \item a main lobe;
    \item side lobes.
\end{itemize}

Tapered windows, such as Hamming or Blackman-Harris windows, reduce side lobes but widen the main lobe.

\begin{authornote}
There is a trade-off. A rectangular window is better when we need to separate close frequency components of similar power. A tapered window is better when leakage from a strong component could hide a weaker component at another frequency.
\end{authornote}

\subsection{Periodogram and Power Spectral Density}

The periodogram is the squared magnitude of the DFT of a time-limited signal. It estimates how signal power is distributed across frequency.

For stochastic signals, a single periodogram can have high variability. Therefore, one often estimates the Power Spectral Density, PSD, using averaging methods.

The PSD is a statistical description of signal power as a function of frequency.

\subsection{Averaged Periodogram and Welch Method}

The averaged periodogram method divides a long signal into shorter epochs. The spectrum of each epoch is computed and then spectra are averaged bin by bin.

Welch's method improves this approach by:

\begin{itemize}[noitemsep]
    \item splitting the signal into possibly overlapping epochs;
    \item applying a windowing function to each epoch;
    \item computing the spectrum of each windowed epoch;
    \item averaging spectra bin by bin.
\end{itemize}

This reduces the variance of the PSD estimate, but at the cost of lower frequency resolution.

\begin{authornote}
Welch's method is useful because biological signals are noisy. Averaging spectra makes the estimate more stable, but shorter windows mean worse frequency resolution.
\end{authornote}

\subsection{Short-Time Fourier Transform and Spectrogram}

The Short-Time Fourier Transform, STFT, estimates the spectrum of a signal in short moving time windows.

It is used to analyze non-stationary signals whose frequency content changes over time.

The spectrogram displays:

\begin{itemize}[noitemsep]
    \item time on the horizontal axis;
    \item frequency on the vertical axis;
    \item magnitude or power using color.
\end{itemize}

\begin{authornote}
STFT introduces the classical time-frequency trade-off. Short windows give good time resolution but poor frequency resolution. Long windows give good frequency resolution but poor time resolution.
\end{authornote}

\subsection{Digital Filters}

A digital filter allows some spectral components of a signal to pass almost unaltered while attenuating other components.

The four basic types are:

\begin{itemize}[noitemsep]
    \item low-pass filter;
    \item high-pass filter;
    \item band-pass filter;
    \item band-stop or band-reject filter.
\end{itemize}

A notch filter is a narrow band-stop filter, commonly used to attenuate power-line noise at 50 Hz, or 60 Hz in some countries.

\subsection{Passband, Stopband, Transition Band, and Cutoff}

The passband is the frequency interval in which the filter gain is close to 1.

The stopband is the frequency interval in which the filter gain is close to 0.

The transition band is the region between passband and stopband.

The cutoff frequency, or corner frequency, marks the limit of the passband. In many filter designs, the gain at cutoff is approximately:

\[
\frac{1}{\sqrt{2}}
\approx
0.71,
\]

corresponding to:

\[
-3\,\text{dB}.
\]

The roll-off is the slope of the frequency response in the transition band. A fast roll-off means a narrow transition band.

\subsection{Good Filter Properties}

Good filter properties include:

\begin{itemize}[noitemsep]
    \item fast roll-off;
    \item flat passband, with no or little ripple;
    \item strong stopband attenuation;
    \item minimal distortion of relevant signal components.
\end{itemize}

Stopband attenuation is better evaluated in decibels rather than on a linear gain scale, because small gains such as $0.001$ and $0.0001$ are difficult to distinguish visually in linear scale.

\subsection{FIR and IIR Filters}

Digital filters are classified into:

\begin{itemize}[noitemsep]
    \item FIR filters, Finite Impulse Response;
    \item IIR filters, Infinite Impulse Response.
\end{itemize}

The output of an FIR filter is a linear combination of input samples.

The output of an IIR filter depends on both input samples and previous output samples.

IIR filters are usually more efficient than FIR filters, meaning that a lower order is often sufficient to reach similar frequency-domain specifications.

However, FIR filters can be designed to have linear phase, while IIR filters generally cannot.

\begin{authornote}
Linear phase means that all frequency components are delayed by the same amount. This is important when preserving waveform shape matters, for example in ERP analysis.
\end{authornote}

\begin{noexam}
The additional filter draft material provides deeper examples of convolution, nonlinear phase effects, and extra demonstrations. These are useful for intuition, but they should be treated as supplementary unless explicitly included in the exam program.
\end{noexam}

\subsection{Core Exam Takeaways}

\begin{enumerate}
    \item ADC consists of sampling and quantization.
    \item Sampling discretizes time; quantization discretizes amplitude.
    \item The Nyquist frequency is half of the sampling frequency.
    \item Aliasing occurs when the signal contains components above Nyquist before sampling.
    \item Aliasing must be prevented using analog anti-aliasing filtering before ADC.
    \item Quantization noise depends on the quantization interval, LSB.
    \item A large ADC range increases quantization error; a small range increases clipping risk.
    \item Averaging $N$ independent noise realizations reduces standard deviation by $\sqrt{N}$.
    \item The DFT represents a sampled signal in the frequency domain.
    \item For real signals, the spectrum is symmetric and the second half is redundant.
    \item Frequency resolution is determined by the total observation time.
    \item Zero-padding improves visual interpolation but does not add information.
    \item Windowing controls spectral leakage.
    \item Welch's method reduces PSD variance at the cost of spectral resolution.
    \item Digital filters are low-pass, high-pass, band-pass, or band-stop.
    \item FIR filters can have linear phase; IIR filters are usually more efficient.
\end{enumerate}

\subsection{True Statements}

\subsubsection{ADC, Sampling, Aliasing, and Quantization}

\begin{enumerate}[label=\arabic*.]
    \item Shannon's theorem, or sampling theorem, states that a continuous signal can be properly sampled only if it does not contain frequency components above one half of the sampling rate.
    
    \item In Analog-to-Digital Conversion, ADC, the Nyquist frequency equals half of the sampling frequency.
    
    \item The reconstruction of an analog signal from its sampled version is equivalent to the sum of a set of $\mathrm{sinc}()$ functions, one for each sample, each centered on the time of the respective sample, whose amplitude equals the sample value.
    
    \item Aliasing occurs when an analog signal is sampled outside the conditions set by Shannon's theorem.
    
    \item When a sinusoidal component with frequency $f_0 \in (f_{\text{Nyquist}}, f_{\text{sampling}})$ is sampled at frequency $f_{\text{sampling}}$, aliasing occurs and the reconstructed signal is a sinusoidal component at:
    \[
    f_{\text{aliasing}}
    =
    f_{\text{sampling}} - f_0.
    \]
    
    \item Aliasing can be prevented by applying an analog low-pass filter with cutoff frequency lower than $f_{\text{Nyquist}}$ to the analog signal before it is converted.
    
    \item Quantization is the approximation of the analog value of a sample to the nearest among the allowed quantization levels.
    
    \item Quantization introduces noise whose amplitude is proportional to the width of the quantization interval:
    \[
    \sigma_{\text{noise}} =
    \frac{1}{\sqrt{12}}LSB.
    \]
    
    \item Quantization divides the input range of the ADC into approximately $2^{N_{\text{bits}}}$ intervals, where $N_{\text{bits}}$ is the number of bits of the ADC.
    
    \item Given a fixed number of bits of the ADC, choosing a large input range increases the quantization error, while choosing a small input range increases the chance that the signal is clipped.
    
    \item Appropriate application of an analog filter before conversion may prevent saturation by removing high-amplitude artifacts in specific frequency bands.
\end{enumerate}

\subsubsection{Statistics, Probability, and Noise}

\begin{enumerate}[label=\arabic*.]
    \item A signal can be entirely deterministic, or it can be known only by means of its statistical properties.
    
    \item Probability theory deals with the mathematical modeling of random variables and stochastic processes.
    
    \item Statistics deals with the description of empirical observations and with the estimation of usually unknown parameters of mathematical models.
    
    \item In statistics, a sample is a collection of observations drawn from a population.
    
    \item In signal processing, a sample is a single measurement of a signal at a specific time point.
    
    \item The Average Rectified Value, ARV, is obtained by summing the absolute values of all samples and dividing by the number of samples.
    
    \item The Root Mean Square, RMS, is the square root of the average of the squared values of the samples of a signal.
    
    \item The variance of a signal is estimated by summing the squared deviations of the sample values from the sample mean and dividing by $N-1$.
    
    \item The standard deviation of a signal is the square root of its variance.
    
    \item In white noise, all samples are uncorrelated.
    
    \item The frequency spectrum of white noise is flat.
    
    \item In Gaussian noise, the probability density that a sample has a given amplitude follows the normal distribution, usually with zero mean.
    
    \item The Central Limit Theorem states that the probability distribution of the average of $N$ independent and identically distributed random variables tends to a normal distribution for $N$ approaching infinity.
    
    \item The probability distribution of the average of $N$ independent and identically distributed Gaussian random variables is normal independently of the value of $N$.
    
    \item Given $N$ independent and identically distributed random variables with variance $\sigma^2$, the variance of their average is $\sigma^2/N$.
    
    \item For independent and identically distributed random variables with standard deviation $\sigma$, the standard deviation of their average is $\sigma/\sqrt{N}$.
    
    \item The synchronized average of $N$ trials containing only spontaneous EEG with RMS equal to $\sigma$ has RMS equal to $\sigma/\sqrt{N}$.
\end{enumerate}

\subsubsection{Fourier Transform, DFT, and Spectral Analysis}

\begin{enumerate}[label=\arabic*.]
    \item Signals can be decomposed into sums of sine waves, each carrying power at a single frequency.
    
    \item The Discrete Fourier Transform associates the time-limited and sampled time-domain representation of a signal with its bandwidth-limited and sampled frequency-domain representation.
    
    \item Given a sampled signal with $N$ samples, its DFT contains $N$ complex values.
    
    \item For a real time-domain signal, the first half of the DFT is sufficient to completely define the signal.
    
    \item The DFT represents the amplitude and initial phase of sine wave components of the signal.
    
    \item Only components up to the Nyquist frequency are needed to characterize the signal, since the remaining part of the spectrum is a mirrored copy of the first half.
    
    \item Fast Fourier Transform, FFT, is an efficient implementation of the DFT, especially for signals whose number of samples is a power of two.
    
    \item The spectrum of a signal is often represented as magnitude and phase plots.
    
    \item Spectral analysis is well suited to identifying narrowband useful signals in approximately white noise.
    
    \item For a signal sampled with sampling interval $\Delta T_s = 1/f_s$, the spectrum has a frequency range of width $f_s$.
    
    \item For a signal of total duration $T_{\max}$, adjacent DFT frequency bins are $\Delta f = 1/T_{\max}$ apart.
    
    \item Zero-padding a signal reduces the distance between adjacent displayed samples of the spectrum, visually improving frequency resolution.
    
    \item Spectral leakage is observed when comparing the spectrum of a signal with the spectrum of a short section of the same signal.
    
    \item Extracting a section from a longer signal is equivalent to multiplying the signal by a rectangular windowing function.
    
    \item Spectral leakage spreads each frequency sample of the original spectrum with a pattern defined by the DFT of the windowing function.
    
    \item The effects of spectral leakage can be controlled by using a tapered windowing function instead of a rectangular one.
    
    \item A rectangular window is preferable when spectral analysis aims to distinguish two spectral components with similar frequency and power.
    
    \item A tapered window is preferable when leakage from a strong component may obscure a much weaker component at a different frequency.
    
    \item The spectrum of a time-limited sampled signal is formally defined as the squared absolute value of its DFT, called periodogram.
    
    \item In the averaged periodogram method, a long signal is split into shorter adjacent epochs, the spectrum of each epoch is computed, and these spectra are averaged bin by bin.
    
    \item In Welch's method, the signal is split into shorter possibly overlapping epochs, each epoch is multiplied by a windowing function, and spectra are averaged bin by bin.
    
    \item The Short-Time Fourier Transform is a simple method to estimate a spectrogram.
    
    \item In a spectrogram, time is usually displayed on the horizontal axis, frequency on the vertical axis, and magnitude or power is color-coded.
\end{enumerate}

\subsubsection{Digital Filters}

\begin{enumerate}[label=\arabic*.]
    \item The purpose of a filter is to allow some spectral components of a signal to pass almost unaltered while blocking other spectral components.
    
    \item Filters are categorized into four basic types: low-pass, high-pass, band-pass, and band-reject or band-stop.
    
    \item The passband is the interval of frequencies in which the gain of the filter is close to 1.
    
    \item The stopband is the interval of frequencies in which the gain is close to 0.
    
    \item The transition band is the interval in which the gain has intermediate values between 0 and 1.
    
    \item The roll-off of a filter is the slope of its frequency response in the transition band.
    
    \item Roll-off is high when the transition band is narrow.
    
    \item The cutoff frequency, or corner frequency, designates the limit of the passband.
    
    \item In many filter designs, the gain at the cutoff frequency is approximately $0.71$, corresponding to $-3\,\text{dB}$.
    
    \item The gain in the passband can decrease monotonically below 1 when the frequency approaches the cutoff frequency, or it may ripple above and below 1.
    
    \item The frequency response of a filter in the stopband should preferably be evaluated with the gain expressed in dB.
    
    \item Good features of a filter include fast roll-off, flat passband, and strong stopband attenuation.
    
    \item Digital filters are categorized into FIR and IIR filters.
    
    \item The output of FIR filters is a linear combination of input samples.
    
    \item The output of IIR filters combines samples of the input and past samples of the output.
    
    \item The Butterworth filter is a design method in the family of IIR filters.
    
    \item The corner frequency of a high-pass filter is called low cutoff frequency.
    
    \item The corner frequency of a low-pass filter is called high cutoff frequency.
    
    \item Band-pass and band-stop filters have both a low cutoff frequency and a high cutoff frequency.
    
    \item A notch filter is a type of band-stop filter that attenuates the input in a narrow band around the notch frequency.
    
    \item The most common use of a notch filter is to remove the power-line artifact at 50 Hz and its harmonics, or 60 Hz in some countries.
    
    \item An IIR filter is more efficient than a FIR filter, meaning that a FIR filter needs to be of higher order to achieve the same quality specifications.
    
    \item A FIR filter can be designed to have linear phase, meaning it does not introduce time-domain distortions in the waveform of the output signal.
    
    \item IIR filters cannot have linear phase.
\end{enumerate}

\section{Brain-Computer Interfaces}

\subsection{Definition of Brain-Computer Interface}

A Brain-Computer Interface, BCI, is a system that measures brain activity and converts it into an artificial output. This output replaces, restores, enhances, supplements, or improves natural Central Nervous System output.

A key feature of a BCI is that it bypasses the normal neuromuscular pathways. Instead of using nerves and muscles to control the environment, the system extracts information directly from brain activity and translates it into commands.

\[
\text{brain activity}
\rightarrow
\text{feature extraction}
\rightarrow
\text{translation algorithm}
\rightarrow
\text{command}
\rightarrow
\text{feedback}.
\]

\begin{authornote}
A BCI is not simply a classifier applied to EEG. It is a closed-loop control system in which the user learns from feedback and modifies brain activity to improve control.
\end{authornote}

\subsection{BCI Closed Loop}

The BCI operates as a closed-loop pipeline:

\begin{enumerate}[noitemsep]
    \item the user generates an intention or cognitive state;
    \item brain signals are acquired;
    \item relevant features are extracted;
    \item features are translated into a command;
    \item the command acts on the environment;
    \item the user receives feedback.
\end{enumerate}

Feedback is essential because controlling the environment without muscles is unnatural. The brain needs feedback to learn how to modulate its signals effectively.

\begin{authornote}
The feedback loop is what allows BCI learning. Without feedback, the user cannot know whether the attempted mental strategy is producing the desired output.
\end{authornote}

\subsection{Invasive and Non-Invasive BCIs}

BCI signals can be acquired using invasive or non-invasive methods.

Invasive methods include:

\begin{itemize}[noitemsep]
    \item microelectrodes;
    \item single-unit activity;
    \item multi-unit activity;
    \item local field potentials;
    \item electrocorticography.
\end{itemize}

Non-invasive methods include:

\begin{itemize}[noitemsep]
    \item EEG;
    \item NIRS;
    \item fMRI.
\end{itemize}

EEG is one of the most widely used methods for non-invasive BCIs because it is relatively cheap, portable, and has good temporal resolution.

\subsection{Assistive and Rehabilitative BCIs}

Clinically, BCIs can be divided into two main paradigms:

\begin{itemize}[noitemsep]
    \item assistive BCIs;
    \item rehabilitative BCIs.
\end{itemize}

Assistive BCIs provide an alternative access channel for users with severe motor impairments. Their goal is not to heal the biological damage, but to allow communication and interaction with the environment.

Rehabilitative BCIs aim to drive neuroplasticity. Their goal is to improve biological function by linking motor intention to contingent sensory feedback.

\begin{authornote}
Assistive BCI replaces an output channel. Rehabilitative BCI tries to retrain the nervous system.
\end{authornote}

\subsection{Assistive Technology and the Accessibility Gap}

Traditional assistive devices, such as joysticks, touch screens, switches, head trackers, and eye trackers, require some reliable motor output.

In conditions such as ALS, locked-in syndrome, or severe stroke, even minimal muscular control may be unavailable. In these cases, a BCI can provide an alternative input channel that bypasses muscles.

The BCI does not need to replace the whole assistive software environment. In a modern modular architecture, the BCI simply outputs a discrete selection command, while standard assistive technology software manages communication, browsing, or domotic control.

\subsection{P300 and the Oddball Paradigm}

The P300 is an Event-Related Potential that appears as a positive voltage deflection approximately 300 ms after the presentation of a rare and meaningful target stimulus.

It is commonly elicited using the oddball paradigm, where rare target stimuli are presented among frequent non-target stimuli.

In a P300 BCI, the user attends to a target item. When that item flashes, a P300 response is generated. When non-target items flash, the P300 is absent or weaker.

\subsection{P300 Speller}

A classical assistive BCI application is the P300 Speller.

The interface consists of a matrix of letters and symbols. Rows and columns flash randomly. The user focuses attention on the desired character.

A P300 is generated when the row or column containing the desired character flashes. By detecting which row and column elicited the P300, the system infers the selected character.

\begin{authornote}
The P300 Speller is not based on motor imagery. It is based on attention to rare target stimuli and ERP detection.
\end{authornote}

\subsection{P300 Feature Extraction and Classification}

From a machine learning perspective, P300 decoding is a binary classification problem:

\[
\text{target}
\quad \text{vs.} \quad
\text{non-target}.
\]

Features are often constructed by taking EEG voltage values at specific post-stimulus latencies and channels. These values may be averaged over short temporal windows, for example around 50 ms.

The resulting feature vector is classified using linear methods such as SWLDA, Stepwise Linear Discriminant Analysis.

\begin{noexam}
Some detailed technical material about P300 feature extraction and SWLDA was marked as not discussed in class in the handout. It is still useful to understand why P300 BCIs work: the classifier is not reading ``letters'' directly, but detecting whether a target-related ERP was present after a flash.
\end{noexam}

\subsection{Sensorimotor Rhythms and Motor Imagery BCIs}

Sensorimotor rhythms, SMR, include rhythms in the mu and beta bands. They are generated over sensorimotor cortical areas.

During actual movement and kinesthetic motor imagery, SMR power decreases. This decrease is often described as Event-Related Desynchronization, ERD.

Motor imagery BCIs exploit the voluntary modulation of SMR power. For example, imagining movement of the left or right hand can modulate rhythms over contralateral sensorimotor areas.

\begin{authornote}
Kinesthetic motor imagery means imagining the feeling of performing a movement, not simply visualizing yourself from outside. This distinction matters because kinesthetic imagery engages motor networks more strongly.
\end{authornote}

\subsection{Asynchronous SMR-BCIs}

Unlike P300 BCIs, which detect discrete event-related responses after stimuli, asynchronous SMR-BCIs translate continuous, state-dependent modulation of brain rhythms into control commands.

For example, in a cursor-control task, the system continuously estimates PSD features in specific EEG bands and channels. These features are translated into a continuous output controlling cursor movement.

\begin{authornote}
P300 BCI is event-driven: the system asks, ``was there a P300 after this flash?'' SMR BCI is state-driven: the system asks, ``what is the current sensorimotor rhythm state?''
\end{authornote}

\subsection{PSD Estimation in Online BCI}

In SMR-BCIs, the system must estimate spectral power in real time. This often uses PSD estimation methods such as Welch's method.

A key design problem is balancing:

\begin{itemize}[noitemsep]
    \item spectral resolution;
    \item update speed;
    \item feedback smoothness;
    \item classification accuracy.
\end{itemize}

The epoch used for spectral estimation may be longer than the update interval. Therefore, consecutive epochs overlap.

This system-level epoch overlap is different from the internal window overlap used by Welch's method.

\begin{noexam}
This technical distinction was marked as additional material, but it is important for understanding real-time BCI design. The BCI can update every 100 ms while using a longer 300 ms data epoch, because consecutive epochs overlap.
\end{noexam}

\subsection{Rehabilitative BCI and Sensorimotor Loop}

After stroke, the connection between motor intention and movement-related sensory feedback may be disrupted. A rehabilitative BCI can act as a synthetic bridge.

The patient performs kinesthetic motor imagery. If the BCI detects appropriate motor-related brain activity, it triggers an external device, such as Functional Electrical Stimulation, FES, or a robotic orthosis.

This generates sensory feedback contingent on motor intent, closing the sensorimotor loop.

\[
\text{motor intent}
\rightarrow
\text{BCI detection}
\rightarrow
\text{external movement or stimulation}
\rightarrow
\text{sensory feedback}.
\]

\begin{authornote}
The rehabilitative idea is Hebbian: pairing motor cortical activation with the correct sensory feedback may reinforce surviving pathways and promote useful plasticity.
\end{authornote}

\subsection{Hybrid BCIs and Corticomuscular Coherence}

Hybrid BCIs combine multiple signal sources. A hybrid system may use both EEG and EMG.

Corticomuscular Coherence, CMC, measures functional synchronization between cortical motor activity and peripheral muscle activation.

A CMC-based rehabilitative BCI can encourage correct coupling between motor cortex and target muscles, while discouraging compensatory strategies such as shoulder activation or mirror movements of the healthy limb.

\begin{noexam}
The handout marks the advanced CMC architecture as optional cutting-edge research. It is useful to understand the concept, but the detailed spectral mathematics behind coherence is beyond the core requirements.
\end{noexam}

\subsection{Clinical Relevance and Brain Network Reconfiguration}

BCI-supported rehabilitation can induce training-specific brain network reconfiguration. In the reported clinical framework, BCI training was associated with improved clinical outcomes and increased involvement of ipsilesional sensorimotor areas.

Interhemispheric connectivity changes can also be band-specific. In particular, increased coupling in beta bands is especially relevant because beta rhythms are directly related to motor function.

\begin{authornote}
This is an important translational point: a successful rehabilitative BCI should not only improve a score, but also produce measurable neurophysiological changes consistent with motor recovery.
\end{authornote}

\subsection{Core Exam Takeaways}

\begin{enumerate}
    \item A BCI measures CNS activity and converts it into an artificial output.
    \item A BCI bypasses normal neuromuscular pathways.
    \item BCI operation is closed-loop and depends on feedback.
    \item Assistive BCIs provide alternative access channels.
    \item Rehabilitative BCIs aim to promote neuroplasticity and restore function.
    \item P300 BCIs exploit rare target-related ERP responses.
    \item The P300 Speller detects which row and column elicit a P300.
    \item SMR-BCIs exploit voluntary modulation of mu and beta rhythms.
    \item Motor imagery produces SMR desynchronization.
    \item Online SMR-BCIs require real-time PSD estimation.
    \item Rehabilitative BCIs close a broken sensorimotor loop using contingent feedback.
    \item Hybrid CMC-BCIs combine EEG and EMG to encourage correct corticomuscular coupling.
\end{enumerate}

\subsection{True Statements}

\begin{enumerate}[label=\arabic*.]
    \item A BCI is a system that measures brain activity and converts it to artificial output that replaces, restores, enhances, supplements, or improves natural brain output and thereby changes the ongoing interactions between the brain and the environment.
    
    \item The P300 ERP generated by attending a target stimulus is exploited to build virtual keyboards based on a BCI.
    
    \item The amplitude of sensorimotor rhythms can be voluntarily modulated through motor imagery to build cursor control based on a BCI.
    
    \item A Brain-Computer Interface completely bypasses normal neuromuscular pathways to measure Central Nervous System activity and convert it into artificial output.
    
    \item The operation of a Brain-Computer Interface is a closed-loop pipeline that relies on environmental feedback to allow the user's brain to learn and consciously modify its signals.
    
    \item In assistive applications, BCIs do not aim to heal biological damage, but rather serve as permanent, alternative access channels for users with severe motor impairments.
    
    \item The P300 is an Event-Related Potential that manifests as a positive voltage deflection approximately 300 milliseconds after the presentation of a rare, meaningful target stimulus.
    
    \item In a P300 BCI, single features are typically constructed by averaging EEG voltage over a short temporal window at specific post-stimulus latencies and spatial channels.
    
    \item Linear classifiers like SWLDA are preferred in P300 BCIs because they generalize well even with small training sets and have interpretable weights.
    
    \item In a modern modular Assistive Technology architecture, a P300 BCI solely outputs a discrete selection trigger, delegating complex application logic to mainstream commercial software.
    
    \item Sensorimotor rhythms, SMR, include the mu rhythm and have spectral components in the alpha and beta bands.
    
    \item Sensorimotor rhythms desynchronize and decrease in spectral power during both physical execution and kinesthetic imagination of movement.
    
    \item Rather than detecting discrete events, asynchronous rehabilitative BCIs translate continuous, state-dependent modulation of SMR power into active motor intent commands.
    
    \item To ensure a smooth and responsive cursor, the BCI system's update interval is set shorter than the minimum epoch duration required for an acceptable spectral estimate.
    
    \item System-level overlap of consecutive data epochs is distinct from the algorithmic overlap of the shorter windows used internally by Welch's method.
    
    \item Although Event-Related Desynchronization formally refers to a non-stationary phenomenon anchored to a discrete event, the term ERD is frequently used in BCI contexts to describe continuous, state-dependent drops in SMR power.
    
    \item The primary goal of a rehabilitative BCI is to guide Central Nervous System plasticity by actively engaging damaged motor networks to restore impaired biological functions.
    
    \item Kinesthetic motor imagery involves rehearsing the actual physiological sensation of moving a limb and successfully engages primary and supplementary motor areas.
    
    \item A rehabilitative BCI acts as a synthetic bridge that closes the broken sensorimotor loop by delivering exogenous sensory feedback exactly when the patient's motor cortex generates endogenous motor intent.
    
    \item Successful BCI-supported motor imagery training induces training-specific brain network reconfiguration, significantly increasing interhemispheric connectivity in beta bands.
\end{enumerate}

\section{EMG}

\subsection{Definition and Applications of EMG}

Electromyography, EMG, is a method used to record the electrical activity generated by skeletal muscles.

EMG is used in several fields:

\begin{itemize}[noitemsep]
    \item neuromuscular physiology;
    \item diagnosis of neuromuscular disorders;
    \item posture analysis;
    \item ergonomics and occupational medicine;
    \item musculoskeletal physical therapy;
    \item gait analysis;
    \item sport science;
    \item EMG biofeedback;
    \item man-machine interfaces.
\end{itemize}

\begin{authornote}
EMG measures muscle electrical activity, not force directly. However, under controlled conditions, EMG amplitude is related to the level of muscle activation.
\end{authornote}

\subsection{Muscle Fibers and Excitability}

A muscle fiber is a single cell of muscular tissue. Its membrane is excitable, meaning that it can respond to depolarization by producing an action potential.

The action potential of a muscle fiber is called Muscle Fiber Action Potential, MFAP.

\[
\text{depolarization}
\rightarrow
\text{MFAP}.
\]

MFAPs propagate along the muscle fiber and are the elementary electrical sources contributing to the EMG signal.

\subsection{Muscle Fiber Conduction Velocity}

The speed at which an MFAP travels along the muscle fiber is called Muscle Fiber Conduction Velocity, MFCV.

Typical values are:

\[
MFCV \approx 2-6\,\text{m/s}.
\]

MFCV depends on several factors:

\begin{itemize}[noitemsep]
    \item ionic concentrations;
    \item intracellular pH;
    \item temperature;
    \item muscle fiber diameter and morphology;
    \item muscle fiber length;
    \item fiber type, such as slow or fast-twitch fibers;
    \item muscle fatigue;
    \item neuromuscular pathology.
\end{itemize}

\begin{authornote}
Fatigue often changes muscle fiber conduction velocity and spectral content of EMG. This is one reason why EMG frequency-domain features can be useful in fatigue analysis.
\end{authornote}

\subsection{Motor Neuron and Motor Unit}

A Motor Unit, MU, is the fundamental functional unit of the neuromuscular system. It consists of:

\begin{itemize}[noitemsep]
    \item one alpha motor neuron;
    \item all the muscle fibers innervated by that motor neuron.
\end{itemize}

When the motor neuron fires, all fibers belonging to that motor unit are activated.

\[
\text{one alpha motor neuron}
+
\text{its muscle fibers}
=
\text{Motor Unit}.
\]

\subsection{Motor Unit Action Potential}

The Motor Unit Action Potential, MUAP, is the electrical signal generated by a single motor unit.

It results from the spatial and temporal summation of the MFAPs generated by the fibers belonging to that motor unit:

\[
MUAP
=
\sum MFAP.
\]

\begin{authornote}
The surface EMG signal is not the activity of a single fiber. It is the superposition of many MUAPs from many active motor units, filtered by body tissues and electrode configuration.
\end{authornote}

\subsection{Force Modulation: Recruitment and Rate Coding}

The Central Nervous System modulates muscle force mainly through two mechanisms:

\begin{itemize}[noitemsep]
    \item recruitment;
    \item rate coding.
\end{itemize}

Recruitment is the activation of additional motor units. More active motor units generally produce greater force.

Rate coding is the increase in firing frequency of already active motor units.

Both mechanisms occur together, but their relative contribution changes with force level.

\begin{itemize}[noitemsep]
    \item at low force levels, recruitment is dominant;
    \item at high force levels and during rapid contractions, rate coding becomes more important.
\end{itemize}

\begin{authornote}
A simple intuition is: to increase force, the CNS can either recruit more motor units or make already active motor units fire faster.
\end{authornote}

\subsection{EMG Amplitude and Muscle Activation}

When more motor units are recruited and/or their firing rate increases, the EMG amplitude tends to increase.

However, the relationship between EMG amplitude and force is not always perfectly linear. It depends on:

\begin{itemize}[noitemsep]
    \item muscle geometry;
    \item electrode placement;
    \item contraction type;
    \item fatigue;
    \item cross-talk;
    \item tissue properties;
    \item motor unit synchronization.
\end{itemize}

\begin{authornote}
EMG amplitude is often used as an estimate of muscle activation, but it should not be interpreted as a direct force sensor without considering experimental conditions.
\end{authornote}

\subsection{Volume Conduction}

The electrical potentials generated by muscle fibers propagate through biological tissues before reaching the recording electrodes.

These tissues include:

\begin{itemize}[noitemsep]
    \item muscle tissue;
    \item fat;
    \item skin;
    \item connective tissue.
\end{itemize}

Together, they act as a volume conductor.

Biological tissues spatially and temporally low-pass filter the EMG signal. This means that high spatial and high temporal frequency components are attenuated before reaching the electrodes.

\begin{authornote}
Volume conduction explains why surface EMG is a blurred version of the underlying muscle electrical activity. The electrode does not see the source directly; it sees the source after propagation through tissues.
\end{authornote}

\subsection{Tripole Model of the MFAP}

The propagating MFAP can be modeled as a tripole, with:

\begin{itemize}[noitemsep]
    \item a central negative phase;
    \item two leading and trailing positive phases.
\end{itemize}

As the tripole propagates along the muscle fiber, the recorded potential changes depending on the relative position of the source and the electrode.

End-fiber effects occur when the action potential reaches the end of the muscle fiber. These components may decay with distance differently from the propagating components.

\subsection{Electrode Types and Placement}

EMG can be recorded using different electrode types:

\begin{itemize}[noitemsep]
    \item surface electrodes;
    \item needle or intramuscular electrodes.
\end{itemize}

Surface EMG is non-invasive and suitable for global muscle activation analysis. Intramuscular EMG is more selective but invasive.

Electrode placement strongly affects the recorded signal. Important factors include:

\begin{itemize}[noitemsep]
    \item electrode position relative to the muscle belly;
    \item distance from the innervation zone;
    \item orientation with respect to muscle fibers;
    \item inter-electrode distance;
    \item skin preparation;
    \item cross-talk from nearby muscles.
\end{itemize}

\begin{authornote}
Bad electrode placement can change EMG amplitude and frequency content more than the physiological effect you are trying to measure. Placement is therefore not a minor detail.
\end{authornote}

\subsection{Selectivity and Cross-Talk}

Selectivity is the ability of the EMG recording to capture activity from the target muscle while rejecting activity from other muscles.

Cross-talk occurs when the EMG signal includes activity from nearby muscles.

Cross-talk is influenced by:

\begin{itemize}[noitemsep]
    \item electrode size;
    \item inter-electrode distance;
    \item depth of the target muscle;
    \item thickness of subcutaneous tissue;
    \item anatomical proximity of muscles.
\end{itemize}

\begin{authornote}
Surface EMG is convenient but not perfectly selective. If two nearby muscles activate together, the electrode may pick up both.
\end{authornote}

\subsection{EMG Signal Processing}

Typical EMG processing includes:

\begin{enumerate}[noitemsep]
    \item acquisition and amplification;
    \item filtering;
    \item artifact inspection;
    \item rectification;
    \item smoothing or envelope extraction;
    \item feature extraction.
\end{enumerate}

A common preprocessing pipeline is:

\[
\text{raw EMG}
\rightarrow
\text{band-pass filtering}
\rightarrow
\text{rectification}
\rightarrow
\text{low-pass smoothing}
\rightarrow
\text{envelope}.
\]

Rectification transforms negative values into positive values:

\[
x_{\text{rect}}(t)
=
|x(t)|.
\]

The envelope represents the slow modulation of EMG amplitude over time.

\begin{authornote}
Rectification is necessary because raw EMG oscillates around zero. Without rectification, positive and negative phases would cancel during averaging or smoothing.
\end{authornote}

\subsection{EMG Features}

EMG features can be extracted in several domains:

\begin{itemize}[noitemsep]
    \item amplitude domain;
    \item time domain;
    \item frequency domain.
\end{itemize}

Common amplitude and time-domain features include:

\begin{itemize}[noitemsep]
    \item RMS;
    \item ARV or MAV, mean absolute value;
    \item integrated EMG;
    \item zero crossings;
    \item waveform length.
\end{itemize}

Frequency-domain features include:

\begin{itemize}[noitemsep]
    \item power spectral density;
    \item mean frequency;
    \item median frequency.
\end{itemize}

\begin{authornote}
RMS and ARV are often used as proxies for muscle activation level. Mean and median frequency are often used when studying fatigue or changes in conduction velocity.
\end{authornote}

\subsection{EMG in Man-Machine Interfaces}

EMG can be used as an input signal for man-machine interfaces. For example, muscle activation patterns can be translated into commands for:

\begin{itemize}[noitemsep]
    \item prosthetic control;
    \item orthotic control;
    \item rehabilitation devices;
    \item gesture recognition systems;
    \item human-computer interaction.
\end{itemize}

Compared with EEG, EMG usually has larger amplitude and higher signal-to-noise ratio, but it depends on the availability of residual muscular activity.

\subsection{Core Exam Takeaways}

\begin{enumerate}
    \item EMG records electrical activity generated by skeletal muscles.
    \item A muscle fiber is an excitable cell.
    \item Depolarization of a muscle fiber produces an MFAP.
    \item MFCV is typically in the range $2$--$6\,\text{m/s}$.
    \item A motor unit consists of one alpha motor neuron and all the muscle fibers it innervates.
    \item A MUAP is the spatial and temporal sum of MFAPs from one motor unit.
    \item Muscle force is modulated by recruitment and rate coding.
    \item Recruitment dominates at low force levels.
    \item Rate coding becomes more important at high forces and during rapid contractions.
    \item Surface EMG is affected by volume conduction.
    \item Biological tissues act as spatial and temporal low-pass filters.
    \item The propagating MFAP can be modeled as a tripole.
    \item Electrode placement strongly affects the recorded EMG.
    \item Cross-talk is contamination from nearby muscles.
    \item EMG processing often includes filtering, rectification, smoothing, and feature extraction.
\end{enumerate}

\subsection{True Statements}

\begin{enumerate}[label=\arabic*.]
    \item A muscle fiber is a single cell of the muscular tissue.
    
    \item The membrane of the muscle fiber is an excitable tissue that produces a Muscle Fiber Action Potential, MFAP, upon depolarization.
    
    \item The speed at which an MFAP travels along the muscle fiber is the Muscle Fiber Conduction Velocity, MFCV.
    
    \item MFCV is typically in the range of $2$--$6\,\text{m/s}$.
    
    \item MFCV is influenced by factors such as muscle temperature, fiber diameter, and fatigue.
    
    \item A Motor Unit, MU, is the fundamental functional unit of the neuromuscular system.
    
    \item A Motor Unit consists of a single alpha motor neuron and all the muscle fibers it innervates.
    
    \item The Motor Unit Action Potential, MUAP, is the electrical signal generated by a single motor unit.
    
    \item The MUAP represents the spatial and temporal sum of all constituent MFAPs.
    
    \item The Central Nervous System modulates muscle force through two primary mechanisms: recruitment and rate coding.
    
    \item Recruitment is the process of activating additional motor units to increase the force of contraction.
    
    \item Rate coding refers to increasing the firing frequency, in impulses per second, of already active motor units.
    
    \item At low force levels, recruitment is the dominant mechanism for force modulation.
    
    \item At high force levels and during rapid contractions, rate coding becomes a more significant contributor to force generation.
    
    \item The combination of recruitment and rate coding differs between increasing and decreasing force.
    
    \item More recruited motor units and greater rate coding generally mean that EMG amplitude increases.
    
    \item The tissues between the muscle fibers and the recording electrodes, such as fat and skin, act as a volume conductor.
    
    \item The volume conductor spatially and temporally low-pass filters the EMG signal.
    
    \item The propagating MFAP can be modeled as a tripole, with a central negative phase and two leading or trailing positive phases.
    
    \item End-fiber effects fall off with distance less than the other components of the MUAP.
    
    \item Temporal high-pass and low-pass components of the MUAP are linked to its spatial components because the source travels with a specific conduction velocity.
\end{enumerate}

